\chapter{Literature Review}
\label{ch:background}

This chapter reviews five works (2019--2023) that formulate a multi-spacecraft coverage
problem and solve it with a meta-heuristic: a \acf{GA}, \acf{PSO}, \ac{NSGA}-II,
\ac{NSGA}-III and \acf{SA}. Table~\ref{tab:lit} summarises their formulations.

\section{Swarm co-orbits for planetary moon mapping}
Nallapu and Thangavelautham~\cite{nallapu2021} design swarms that map the moon Deimos.
The variables mix integers (number of encounters, spacecraft per encounter) with reals
(encounter angles). The objective is $\min N$ subject to coverage $\ge 90\,\%$, no
collisions and a fuel budget. A GA is compared with PSO, which uses a penalty for
infeasible designs. GA found 3 spacecraft for 90.5\,\% coverage, PSO found 4.
\emph{Strength:} a fair GA/PSO comparison. \emph{Shortcoming:} only one objective, and
hidden terrain is not checked.

\section{Swarm flyby mapping of asteroid Eros}
Nallapu and Thangavelautham~\cite{nallapu2019} place $N$ spacecraft on a ring around
Eros. The variables are $N\in[1,20]$ and each spacecraft's angle $\theta_i\in[0,2\pi]$;
the problem is $\min N$ subject to $100\,\%$ surface coverage. A mixed-integer GA found
that 8 spacecraft cover all of Eros in one pass. \emph{Shortcoming:} the spacecraft only point at the centre,
lighting and blocked views are ignored.

\section{Multi-mission CubeSat constellation design}
Melaku and Kim~\cite{melaku2023} optimise a CubeSat constellation using inclination $i$,
number of planes $P$ and satellites per plane $S$. They minimise $PS$ and the average
revisit time, and maximise coverage, with NSGA-II (population 100, 500 generations). For
example, 20 satellites give 98.8\,\% coverage. \emph{Strength:} a Pareto front of
trade-offs. \emph{Shortcoming:} only three variables, and the Earth is a sphere.

\section{Walker constellation design for LEO navigation}
Guan et al.~\cite{guan2020} design navigation constellations with variables $S$, $P$, $i$
and altitude. They minimise the worst-case positioning error (GDOP) and the number of
satellites with NSGA-III, and use a GA for a single-objective case. For example, 180
satellites at 1500\,km give GDOP $<3$. \emph{Strength:} a worst-case objective and parallel
computing for slow evaluations. \emph{Shortcoming:} navigation rather than imaging.

\section{Swarm orbital formation around asteroids}
Stolfi Rosso and Danoy~\cite{stolfi2023} spread spacecraft evenly around an asteroid to
maximise coverage while saving fuel. A GA tunes the formation parameters for 2, 3, 5 and
10 satellites, and SA tunes each spacecraft on board; SA was more accurate over 400
scenarios. \emph{Strength:} shows SA works for this kind of problem. \emph{Shortcoming:}
even spacing is used instead of measured coverage.

\begin{table}[H]
\centering
\caption{Summary of the reviewed formulations}
\label{tab:lit}
\footnotesize
\begin{tabularx}{\textwidth}{@{}p{1.1cm}p{2.9cm}Xp{1.9cm}p{2.7cm}@{}}
\toprule
Work & Decision variables & Objectives & Algorithm & Key result \\
\midrule
\cite{nallapu2021} 2021 & encounters, spacecraft per encounter, angles & $\min N$ s.t.\ coverage, collision, fuel & GA, PSO & 3 S/C, 90.5\,\% \\
\cite{nallapu2019} 2019 & $N$, ring angles $\theta_i$ & $\min N$ s.t.\ 100\,\% coverage & GA & 8 S/C, 100\,\% of Eros \\
\cite{melaku2023} 2023 & $i$, $P$, $S$ & $\min PS$, $\min$ revisit, $\max$ coverage & NSGA-II & 20 S/C, 98.8\,\% \\
\cite{guan2020} 2020 & $S$, $P$, $i$, altitude & $\min$ worst GDOP, $\min SP$ & NSGA-III, GA & 180 S/C, GDOP $<3$ \\
\cite{stolfi2023} 2023 & formation parameters & even spread, fuel & GA, SA & SA best in 400 cases \\
\bottomrule
\end{tabularx}
\end{table}

\textbf{Metrics and gap.} From these works we take coverage as the main objective, the number of
spacecraft and the fuel as costs, and the best, mean and standard deviation of the fitness over
repeated runs to compare algorithms. \emph{Gap:} the reviewed works either fix the trajectories
(a single flyby or a ring), treat the body as convex or fully lit, or optimise the number of
spacecraft separately from their orbits. Our formulation decides the fleet size, every orbit and
the mission duration together, and measures coverage on the real non-convex shape with occlusion,
sunlight, cast shadows and an image-resolution limit. The comparison with
Nallapu and Thangavelautham~\cite{nallapu2019} is repeated with our model in
Section~\ref{sec:v5}.
